% GENERATED FILE. Edit canonical lessons, then run npm run generate:books.
% canonical-source-hash: fnv1a64:f10e84ae0a89c6aa
% canonical-lessons: MR-C246-indradhanusya, MR-C246-gara, MR-C246-vasant, MR-C246-grah, MR-C246-svatantrya

\chapter{Rainbow, Hail, Springtime, Planet, Freedom}
\label{ch:mr-rainbow-hail-springtime-planet-freedom}
\hlchaptermodality{246}

\begin{chapteropening}
\textbf{By the end of this chapter:} \emph{I can say a rainbow, hail, springtime, a planet and freedom.}

Complete the last lesson of chapter 246: I can say a rainbow, hail, springtime, a planet and freedom.
\end{chapteropening}

\section[\texorpdfstring{\mr{इंद्रधनुष्य} (indradhanuṣya)}{इंद्रधनुष्य (indradhanuṣya)}]{\mr{इंद्रधनुष्य} (indradhanuṣya) — a rainbow}

\label{lesson:MR-C246-indradhanusya}



\begin{quote}
Before the new one: say the Marathi for a college, then the Marathi for a museum.
\end{quote}

\subsection*{You'll want to know: \mr{इंद्रधनुष्य}}
\textbf{\mr{इंद्रधनुष्य}} — \emph{indradhanuṣya} — \textquotedblleft{}a rainbow\textquotedblright{}.

\begin{cousinweb}
\textbf{\mr{इंद्रधनुष्य}} (\emph{indradhanuṣya}) is Sanskrit \emph{indradhanus}, \textquotedblleft{}Indra's bow\textquotedblright{}, a tatsama: \emph{Indra}, the sky god, and \emph{dhanus}, \textquotedblleft{}bow\textquotedblright{}.
\end{cousinweb}

\begin{grammarlens}[title={What you've built}]
One more word for this chapter.
\end{grammarlens}

\subsection*{Your turn}
\begin{itemize}
\raggedright
  \item \emph{Say it:} \emph{indradhanuṣya}
  \item \emph{Say it:} \emph{indradhanuṣya}, once more
  \item \emph{Say it:} say \emph{indradhanuṣya}
  \item \emph{Recall:} say the Marathi for a college, then the Marathi for a museum, then say \emph{indradhanuṣya} again
\end{itemize}

\begin{tcolorbox}[breakable,colback=teal!4,colframe=teal!35!black,title={Before you move on}]
What does \mr{इंद्रधनुष्य} mean? (\textquotedblleft{}A rainbow\textquotedblright{}.) Say it once more.
\end{tcolorbox}
\section[\texorpdfstring{\mr{गारा} (gārā)}{गारा (gārā)}]{\mr{गारा} (gārā) — hail}

\label{lesson:MR-C246-gara}



\begin{quote}
Before the new one: say the Marathi for a museum, then the Marathi for a rainbow.
\end{quote}

\subsection*{You'll want to know: \mr{गारा}}
\textbf{\mr{गारा}} — \emph{gārā} — \textquotedblleft{}hail\textquotedblright{}.

\begin{cousinweb}
\textbf{\mr{गारा}} (\emph{gārā}) is a tadbhava, a word inherited from Sanskrit and reshaped by sound change. It comes from \emph{karaka}, \textquotedblleft{}hail\textquotedblright{} (CDIAL 2782). CDIAL lists the Marathi word here with a query.
\end{cousinweb}

\begin{grammarlens}[title={What you've built}]
One more word for this chapter.
\end{grammarlens}

\subsection*{Your turn}
\begin{itemize}
\raggedright
  \item \emph{Say it:} \emph{gārā}
  \item \emph{Say it:} \emph{gārā}, once more
  \item \emph{Say it:} say \emph{gārā}
  \item \emph{Recall:} say the Marathi for a museum, then the Marathi for a rainbow, then say \emph{gārā} again
\end{itemize}

\begin{tcolorbox}[breakable,colback=teal!4,colframe=teal!35!black,title={Before you move on}]
What does \mr{गारा} mean? (\textquotedblleft{}Hail\textquotedblright{}.) Say it once more.
\end{tcolorbox}
\section[\texorpdfstring{\mr{वसंत} (vasant)}{वसंत (vasant)}]{\mr{वसंत} (vasant) — springtime}

\label{lesson:MR-C246-vasant}



\begin{quote}
Before the new one: say the Marathi for a rainbow, then the Marathi for hail.
\end{quote}

\subsection*{You'll want to know: \mr{वसंत}}
\textbf{\mr{वसंत}} — \emph{vasant} — \textquotedblleft{}springtime\textquotedblright{}.

\begin{cousinweb}
\textbf{\mr{वसंत}} (\emph{vasant}) is Sanskrit \emph{vasanta}, \textquotedblleft{}spring season\textquotedblright{}, taken over whole as a learned loan, a tatsama (CDIAL 11439).
\end{cousinweb}

\begin{grammarlens}[title={What you've built}]
One more word for this chapter.
\end{grammarlens}

\subsection*{Your turn}
\begin{itemize}
\raggedright
  \item \emph{Say it:} \emph{vasant}
  \item \emph{Say it:} \emph{vasant}, once more
  \item \emph{Say it:} say \emph{vasant}
  \item \emph{Recall:} say the Marathi for a rainbow, then the Marathi for hail, then say \emph{vasant} again
\end{itemize}

\begin{tcolorbox}[breakable,colback=teal!4,colframe=teal!35!black,title={Before you move on}]
What does \mr{वसंत} mean? (\textquotedblleft{}Springtime\textquotedblright{}.) Say it once more.
\end{tcolorbox}
\section[\texorpdfstring{\mr{ग्रह} (grah)}{ग्रह (grah)}]{\mr{ग्रह} (grah) — a planet}

\label{lesson:MR-C246-grah}



\begin{quote}
Before the new one: say the Marathi for hail, then the Marathi for springtime.
\end{quote}

\subsection*{You'll want to know: \mr{ग्रह}}
\textbf{\mr{ग्रह}} — \emph{grah} — \textquotedblleft{}a planet\textquotedblright{}.

\begin{cousinweb}
\textbf{\mr{ग्रह}} (\emph{grah}) is Sanskrit \emph{graha}, \textquotedblleft{}seizing\textquotedblright{}, taken over whole as a learned loan, a tatsama (CDIAL 4363). The \textquotedblleft{}seizers\textquotedblright{} were the planets, thought to grip human fate.
\end{cousinweb}

\begin{grammarlens}[title={What you've built}]
One more word for this chapter.
\end{grammarlens}

\subsection*{Your turn}
\begin{itemize}
\raggedright
  \item \emph{Say it:} \emph{grah}
  \item \emph{Say it:} \emph{grah}, once more
  \item \emph{Say it:} say \emph{grah}
  \item \emph{Recall:} say the Marathi for hail, then the Marathi for springtime, then say \emph{grah} again
\end{itemize}

\begin{tcolorbox}[breakable,colback=teal!4,colframe=teal!35!black,title={Before you move on}]
What does \mr{ग्रह} mean? (\textquotedblleft{}A planet\textquotedblright{}.) Say it once more.
\end{tcolorbox}
\section[\texorpdfstring{\mr{स्वातंत्र्य} (svātantrya)}{स्वातंत्र्य (svātantrya)}]{\mr{स्वातंत्र्य} (svātantrya) — freedom}

\label{lesson:MR-C246-svatantrya}



\begin{quote}
Before the new one: say the Marathi for springtime, then the Marathi for a planet.
\end{quote}

\subsection*{You'll want to know: \mr{स्वातंत्र्य}}
\textbf{\mr{स्वातंत्र्य}} — \emph{svātantrya} — \textquotedblleft{}freedom\textquotedblright{}.

\begin{cousinweb}
\textbf{\mr{स्वातंत्र्य}} (\emph{svātantrya}) is Sanskrit \emph{svātantrya}, \textquotedblleft{}independence\textquotedblright{}, a tatsama built on \emph{svatantra}: \emph{sva}, \textquotedblleft{}own\textquotedblright{}, and \emph{tantra}, \textquotedblleft{}rule\textquotedblright{}. CDIAL has no entry for it.
\end{cousinweb}

\begin{grammarlens}[title={What you've built}]
One more word for this chapter.
\end{grammarlens}

\subsection*{Your turn}
\begin{itemize}
\raggedright
  \item \emph{Say it:} \emph{svātantrya}
  \item \emph{Say it:} \emph{svātantrya}, once more
  \item \emph{Say it:} say \emph{svātantrya}
  \item \emph{Recall:} say the Marathi for springtime, then the Marathi for a planet, then say \emph{svātantrya} again
\end{itemize}

\begin{tcolorbox}[breakable,colback=teal!4,colframe=teal!35!black,title={Before you move on}]
What does \mr{स्वातंत्र्य} mean? (\textquotedblleft{}Freedom\textquotedblright{}.) Say it once more.
\end{tcolorbox}
